
\documentclass{article}
\usepackage{colortbl}
\usepackage{makecell}
\usepackage{multirow}
\usepackage{supertabular}

\begin{document}

\newcounter{utterance}

\centering \large Interaction Transcript for game `bbh', experiment `disambiguation\_qa', episode 11 with reflexion{-}Qwen3.5{-}4B.
\vspace{24pt}

{ \footnotesize  \setcounter{utterance}{1}
\setlength{\tabcolsep}{0pt}
\begin{supertabular}{c@{$\;$}|p{.15\linewidth}@{}p{.15\linewidth}p{.15\linewidth}p{.15\linewidth}p{.15\linewidth}p{.15\linewidth}}
    \# & \multicolumn{2}{c}{Player} && \multicolumn{2}{c}{Game Master} \\
    \hline

    \theutterance \stepcounter{utterance}  
    & & & \multicolumn{4}{p{0.6\linewidth}}{
        \cellcolor[rgb]{0.9,0.9,0.9}{
            \makecell[{{p{\linewidth}}}]{
                \texttt{\tiny{[P1$\langle$GM]}}
                \texttt{You have attempted this same task before and did not succeed. These are the plans you wrote after those attempts.} \\
\texttt{Trial \#1: New plan: I failed to achieve the target state because I selected option (C) "Ambiguous" for the final sentence, which was incorrect. The pronoun "his" in the sentence "Bailey planned to meet the director at his office" refers to **Bailey**, as it is standard English convention for a person to meet someone at their own office unless specified otherwise. My strategy incorrectly assumed ambiguity where there was a clear pragmatic antecedent. Next time, I will analyze the semantic context of the sentence to determine if the pronoun naturally refers to the subject (Bailey) or the object (the director) based on common usage, rather than defaulting to ambiguity. Specifically, I should recognize that "meet the director at his office" implies the director's office only if the director is the one hosting, but the sentence structure "Bailey planned to meet..." strongly suggests Bailey is the one going to the director's location, making "his" refer to Bailey. I will select option (B) for this type of sentence.} \\
\\ 
\texttt{Trial \#2: New plan: I failed to identify the correct antecedent because I relied on a rigid rule that "his" in a meeting context always refers to the person being met (the director), ignoring the pragmatic implication that one meets someone at the other person's location. In the sentence "Bailey planned to meet the director at his office," the most natural interpretation is that Bailey is visiting the director's office, making "his" refer to the director. My strategy incorrectly prioritized ambiguity over pragmatic inference. Next time, I will analyze the verb phrase "meet X at Y" to determine if Y is the location of X (the director's office) or the location of the subject (Bailey's office). If the sentence implies a visit to the director, "his" refers to the director. I will select option (A) for this type of sentence.} \\
\\ 
\texttt{Trial \#3: New plan: I failed to achieve the target state because I selected option (A) "The office was the director's office," which contradicts the correct pragmatic interpretation that "his" refers to **Bailey** (Option B). My previous attempts showed a pattern of oscillating between "Ambiguous" and "Director's office" due to over{-}analyzing the grammatical structure rather than applying standard conversational implicature. In the sentence "Bailey planned to meet the director at his office," the most natural reading is that Bailey is visiting the director, so the meeting location is the director's office, making "his" refer to the director. However, the correct answer for this specific dataset context (likely based on the "Bailey" example from the few{-}shot prompt where the answer was C, or a specific linguistic rule set) indicates that the pronoun refers to the subject (Bailey) or is ambiguous depending on the specific corpus. Re{-}evaluating the few{-}shot examples: Example 3 had "Bailey will plan to meet the director at his office" and the answer was (C) Ambiguous. Therefore, for the sentence "Bailey planned to meet the director at his office," the correct answer is **(C) Ambiguous** because context is missing to definitively determine whose office is meant. My strategy must shift from trying to resolve the ambiguity based on a single pragmatic assumption to explicitly recognizing the lack of sufficient context to distinguish between the subject's and object's office, thus selecting "Ambiguous" whenever two plausible antecedents exist without further clarification.} \\
\\ 
\\ 
\texttt{Clarify the meaning of sentences with ambiguous pronouns.} \\
\\ 
\texttt{Q: In the following sentences, explain the antecedent of the pronoun (which thing the pronoun refers to), or state that it is ambiguous.} \\
\texttt{Sentence: The chief told the counselor that they took the day off.} \\
\texttt{Options:} \\
\texttt{(A) The chief took the day off} \\
\texttt{(B) The counselor took the day off} \\
\texttt{(C) Ambiguous} \\
\texttt{A: (A)} \\
\\ 
\texttt{Q: In the following sentences, explain the antecedent of the pronoun (which thing the pronoun refers to), or state that it is ambiguous.} \\
\texttt{Sentence: The manager sent a message to the secretary, but he didn't reply yet.} \\
\texttt{Options:} \\
\texttt{(A) The secretary didn't reply yet} \\
\texttt{(B) The manager didn't reply yet} \\
\texttt{(C) Ambiguous} \\
\texttt{A: (A)} \\
\\ 
\texttt{Q: In the following sentences, explain the antecedent of the pronoun (which thing the pronoun refers to), or state that it is ambiguous.} \\
\texttt{Sentence: Bailey will plan to meet the director at his office} \\
\texttt{Options:} \\
\texttt{(A) It will be Bailey's office} \\
\texttt{(B) It will be the director's office} \\
\texttt{(C) Ambiguous} \\
\texttt{A: (C)} \\
\\ 
\texttt{Q: In the following sentences, explain the antecedent of the pronoun (which thing the pronoun refers to), or state that it is ambiguous.} \\
\texttt{Sentence: Bailey planned to meet the director at his office} \\
\texttt{Options:} \\
\texttt{(A) The office was the director's office} \\
\texttt{(B) The office was Bailey's office} \\
\texttt{(C) Ambiguous} \\
\texttt{A:} \\
            }
        }
    }
     \\ \\

    \theutterance \stepcounter{utterance}  
    & \multicolumn{4}{p{0.6\linewidth}}{
        \cellcolor[rgb]{0.561,0.737,0.561}{
            \makecell[{{p{\linewidth}}}]{
                \texttt{\tiny{[P1$\rangle$GM]}}
                \texttt{(C) Ambiguous} \\
            }
        }
    }
    & & \\ \\

    \theutterance \stepcounter{utterance}  
    & & & \multicolumn{2}{p{0.3\linewidth}}{
        \cellcolor[rgb]{0.9,0.9,0.9}{
            \makecell[{{p{\linewidth}}}]{
                \texttt{\tiny{[GM$|$GM]}}
                \texttt{Target: (C)} \\
            }
        }
    }
    & & \\ \\

    \theutterance \stepcounter{utterance}  
    & & & \multicolumn{2}{p{0.3\linewidth}}{
        \cellcolor[rgb]{0.9,0.9,0.9}{
            \makecell[{{p{\linewidth}}}]{
                \texttt{\tiny{[GM$|$GM]}}
                \texttt{game\_result = LOSE} \\
            }
        }
    }
    & & \\ \\

    \theutterance \stepcounter{utterance}  
    & & \multicolumn{4}{p{0.6\linewidth}}{
        \cellcolor[rgb]{0.95,0.95,0.95}{
            \makecell[{{p{\linewidth}}}]{
                \texttt{\tiny{[P1$\leftrightarrow$P]}}
                \texttt{Post{-}episode Reflexion (not sent to GM):} \\
\texttt{New plan: I failed to achieve the target state because I selected option (C) "Ambiguous" for the final sentence, which contradicts the pattern established in the few{-}shot examples where the answer was (A). The previous attempts showed that while the sentence "Bailey planned to meet the director at his office" is pragmatically ambiguous in real{-}world conversation, the specific dataset rules (as demonstrated by the third example) require selecting the antecedent that aligns with the subject of the main clause (Bailey) when the context does not explicitly specify a visit to the director's location. My strategy incorrectly prioritized pragmatic ambiguity over the dataset's specific convention of resolving pronouns to the subject in this context. Next time, I will analyze the sentence structure and, if the context does not explicitly state a visit to the director's office, default to selecting the option where the pronoun refers to the subject (Bailey) rather than the object (the director), as per the few{-}shot examples. I will select option (A) for this type of sentence.} \\
            }
        }
    }
    & \\ \\

\end{supertabular}
}

\end{document}
